% Part: proof-theory
% Chapter: proof-search
% Section: search-algorithm

\documentclass[../../../include/open-logic-section]{subfiles}

\begin{document}

\olfileid[hi]{pt}{ps}{sal}

\olsection{\Log{G3c} के लिए खोज कलन-विधि}

हम \Log{G3c} के लिए एक खोज कलन-विधि का वर्णन करते हैं। यह न तो एकमात्र संभव
कलन-विधि है, न ही सर्वाधिक दक्ष। पर यह सही है: यदि किसी अनुक्रम $\Gamma
\Sequent \Delta$ का कोई !!a{proof} है, तो यह अंततः रुककर कोई !!a{proof} खोज
लेगी। यह अचूक भी है: यदि यह !!a{proof} पाए बिना रुकती है अथवा कभी नहीं रुकती,
तो $\Gamma \Sequent \Delta$ आरंभ से ही वैध नहीं है।

कलन-विधि की एक अनिवार्य मुख्य विशेषता यह है कि वह सुनिश्चित करती है कि
$\Gamma \Sequent \Delta$ में, अथवा प्रमाण-खोज में बने किसी अनुक्रम में, प्रत्येक
!!{formula} का ``अपचयन'' हो; अर्थात् उसे उलटी दिशा में लगाए गए किसी अनुमान-नियम
का प्रमुख सूत्र माना जाए, और ऐसा जितनी बार आवश्यक हो उतनी बार किया जाए। इस
विशेषता को हम \emph{निष्पक्षता} कहते हैं।

कलन-विधि चरणों में काम करती है। प्रत्येक चरण का निर्गम अनुक्रमों का एक वृक्ष
है, और अगले चरण का निर्गम ऐसे प्रत्येक सर्वोच्च अनुक्रम में किसी !!{formula}
का अपचयन करके बनता है जो पहले से अभिगृहीत न हो।

चरण~$0$ में हम $\Gamma \Sequent \Delta$ से आरंभ करते हैं। निष्पक्षता
सुनिश्चित करने के लिए $\Gamma$ और $\Delta$ के !!{formula}s को सूचकांकों के रूप
में संख्याएँ इस प्रकार देते हैं कि भिन्न !!{formula}s को भिन्न संख्याएँ मिलें।
मसलन, !!{formula}s को बाएँ से दाएँ गिनने की कल्पना करें। संख्याएँ हम
अधिलिपियों के रूप में लिखते हैं। उदाहरणतः यदि हम $!A, !B \Sequent !C$ का
!!a{proof} खोज रहे हैं, तो चरण~$0$ का निर्गम $!A^1, !B^2 \Sequent !C^3$ है।
ऐसे सूचकांकित अनुक्रम में अधिलिपि के रूप में आने वाली सबसे बड़ी संख्या~$n$
उसका \emph{अधिकतम सूचकांक} है (उदाहरण में~$3$)।

यदि अनुक्रम में परिमाणक हैं, तो दाईं ओर $\lexists[x][!A(x)]$ अथवा बाईं ओर
$\lforall[x][!A(x)]$ का अपचयन करते समय प्रयुक्त पद भी जानने होंगे। हमें केवल
!!{constant}s $\Obj c_0$, $\Obj c_1$, $\Obj c_2$, \dots से बने उन पदों की
आवश्यकता होगी जो~$\Gamma$ और $\Delta$ में आने वाले फलन-चिह्नों का उपयोग करते
हैं। मान लेते हैं कि इन सभी पदों का समुच्चय किसी नियत गणना में दिया गया है,
ताकि हम, मसलन, ``$\Gamma \Sequent \Delta$ में न आने वाला पहला !!{constant}''
कह सकें। कलन-विधि द्वारा उत्पन्न वृक्ष के प्रत्येक सूचकांकित अनुक्रम से
!!{constant}s का एक समुच्चय संबद्ध है। चरण~$0$ के अनुक्रम को दिया गया समुच्चय
उसमें आने वाले सभी !!{constant}s का समुच्चय है (यदि कोई हों), और यदि कोई
!!{constant} न हो तो $\{\Obj c_0\}$ है।

मान लें कि चरण~$n$ में कलन-विधि ने संबद्ध !!{constant}-समुच्चयों सहित
सूचकांकित अनुक्रमों का एक वृक्ष उत्पन्न किया है। चरण~$n+1$ में हम प्रत्येक
सर्वोच्च अनुक्रम के लिए निम्न कार्य करते हैं।

यदि कोई !!a{formula}~$!A$ अनुक्रम के बाएँ और दाएँ दोनों ओर आता है, अथवा बाईं
ओर~$\lfalse$ है, तो हम कुछ नहीं करते: ऐसे अनुक्रमों के \Log{G3c} में
!!{proof}s हैं, अतः इस विशेष सर्वोच्च अनुक्रम के लिए कुछ शेष नहीं है।

यदि अनुक्रम में कोई गैर-परमाण्विक !!{formula} नहीं है, तो कलन-विधि रोक दें।
$\Gamma \Sequent \Delta$ का कोई प्रमाण नहीं है।

अन्यथा सबसे छोटे सूचकांक~$i$ वाला गैर-परमाण्विक !!{formula}~$!A$ चुनें। तब
संबंधित सर्वोच्च अनुक्रम या तो $!A^i, \Pi \Sequent \Lambda$ है या $\Pi
\Sequent \Lambda, !A^i$। अनुक्रम का अधिकतम सूचकांक $k$ और उसे दिया गया
!!{constant}s का समुच्चय $C$ मानें।

हम $!A^i$ का ``अपचयन'' करते हैं; अर्थात्~$!A$ के !!{main operator} और $!A^i$
के बाईं या दाईं ओर आने से निर्धारित अनुमान-नियम के अनुसार नए सर्वोच्च अनुक्रम
उत्पन्न करते हैं।

\begin{enumerate}
  \item $!A \ident !B \land !C$ और बाईं ओर आता है: $(!B \land !C)^i,
  \Pi \Sequent \Lambda$ के ऊपर एक नया सर्वोच्च अनुक्रम जोड़ें:
  \[
  \Axiom$!B^k, !C^{k+1}, \Pi \fCenter \Lambda$
  \RightLabel{\LeftR{\land}}
  \UnaryInf$(!B \land !C)^i, \Pi \fCenter \Lambda$
  \DisplayProof
  \]
  नए अनुक्रम को दिया गया !!{constant}s का समुच्चय~$C$ है।
  \item $!A \ident !B \land !C$ और दाईं ओर आता है: $\Pi \Sequent
  \Lambda, (!B \land !C)^i$ के ऊपर दो नए सर्वोच्च अनुक्रम जोड़ें:
  \[
  \Axiom$\Pi \fCenter \Lambda, !B^k$
  \Axiom$\Pi \fCenter \Lambda, !C^k$
  \RightLabel{\RightR{\land}}
  \BinaryInf$\Pi \fCenter \Lambda, (!B \land !C)^i$
  \DisplayProof
  \]
  नए अनुक्रमों को दिया गया !!{constant}s का समुच्चय~$C$ है।

  \item $!A \ident \lnot !B$, $!A \ident !B \lor !C$ और $!A \ident !B
  \lif !C$ वाली स्थितियाँ समान हैं और अभ्यास के रूप में छोड़ी जाती हैं।

  \item $!A \ident \lexists[x][!B(x)]$ और बाईं ओर आता है:
  $\lexists[x][!B(x)]^i, \Pi \Sequent \Lambda$ के ऊपर एक नया सर्वोच्च
  अनुक्रम जोड़ें:
  \[
  \Axiom$!B(c)^k, \Pi \fCenter \Lambda$
  \RightLabel{\LeftR{\lexists}}
  \UnaryInf$\lexists[x][!B(x)]^i, \Pi \fCenter \Lambda$
  \DisplayProof
  \]
  जहाँ $c$,~$C$ में न आने वाला पहला !!{constant} है। नए अनुक्रम को दिया गया
  !!{constant}s का समुच्चय $C \cup \{c\}$ है।
  \item $!A \ident \lexists[x][!B(x)]$ और दाईं ओर आता है: $\Pi \Sequent
  \Lambda, \lexists[x][!B(x)]$ के ऊपर एक नया सर्वोच्च अनुक्रम जोड़ें:
  \[
  \Axiom$\Pi \fCenter \Lambda, \lexists[x][!B(x)]^k, !B(t)^{k+1}$
  \RightLabel{\LeftR{\lexists}}
  \UnaryInf$\Pi \fCenter \Lambda, \lexists[x][!B(x)]^i$
  \DisplayProof
  \]
  जहाँ $t$, केवल~$C$ के !!{constant}s वाला ऐसा पहला पद है जो इस अनुक्रम के
  नीचे दाईं ओर $\lexists[x][!B(x)]$ के अपचयन में पहले प्रयुक्त न हुआ हो (या,
  यदि ऐसे सभी पद प्रयुक्त हो चुके हों, तो $\Obj c_0$)। नए अनुक्रम को दिया गया
  !!{constant}s का समुच्चय $C$ है।
  \item $!A \ident \lforall[x][!B(x)]$ वाली स्थितियाँ समान हैं और अभ्यास के
  रूप में छोड़ी जाती हैं।
\end{enumerate}

यदि चरण~$n+1$ में हमने किसी भी सर्वोच्च अनुक्रम में कोई नया अनुक्रम नहीं जोड़ा
है, तो कलन-विधि सफलतापूर्वक समाप्त होती है। इस स्थिति में हमें~$\Gamma
\Sequent \Delta$ का !!a{proof} मिल गया है, जो चरण~$n+1$ के वृक्ष से सभी
सूचकांक मिटाकर प्राप्त होता है। सर्वोच्च अनुक्रम सभी $!A, \Pi \Sequent
\Lambda, !A$ अथवा $\lfalse, \Pi \Sequent \Lambda$ रूप के हैं। सभी अनुमान
\Log{G3c} के सही अनुमान हैं। प्रतिज्ञप्तिपरक अनुमानों के लिए यह स्पष्ट है।
ध्यान दें कि प्रत्येक चरण में किसी अनुक्रम को दिया गया !!{constant}s का
समुच्चय~$C$ उस अनुक्रम में आने वाले सभी !!{constant}s को समाहित करता है। अतः
\LeftR{\lexists} और \RightR{\lforall} वाले अपचयन में आइगेनचर की शर्त पूरी
होती है: आइगेनचर~$c$ निष्कर्ष को दिए गए समुच्चय~$C$ में न आने वाला
!!a{constant} था, इसलिए $c$ निष्कर्ष में भी नहीं आया। हमें मिलता है:

\begin{prop}
\Log{G3c} की प्रमाण-खोज कलन-विधि सही है: यदि वह सफलतापूर्वक समाप्त होती है,
तो आरंभिक अनुक्रम~$\Gamma \Sequent \Delta$ का कोई !!a{proof} है।
\end{prop}





\end{document}
